\section{Exploratorní analýza. Redukce rozměru dat, hlavní komponenty, faktorová, korespondenční a diskriminační analýza}

\subsection*{Základní myšlenka exploratorní analýzy}

Exploratorní analýza dat, zkráceně EDA, slouží k tomu, abychom data nejdříve pochopili, než na ně
začneme aplikovat formální modely. Cílem není hned dokazovat hypotézy, ale:
\begin{itemize}
    \item zobrazit a shrnout data,
    \item najít zajímavé vzory a struktury,
    \item odhalit odlehlá pozorování a chyby,
    \item ověřit vhodnost proměnných pro další analýzu,
    \item zjednodušit data, pokud mají mnoho proměnných nebo pozorování.
\end{itemize}

EDA typicky používá:
\begin{itemize}
    \item deskriptivní statistiky,
    \item grafy,
    \item korelační analýzu,
    \item redukci dimenze,
    \item shlukování,
    \item metody pro detekci odlehlých hodnot.
\end{itemize}

\subsection*{Příprava dat před redukčními metodami}

Před metodami jako PCA, faktorová analýza, korespondenční analýza nebo diskriminační analýza je potřeba
data zkontrolovat.

\begin{itemize}
    \item \textbf{Deskriptivní statistiky:} průměr, medián, směrodatná odchylka, minimum, maximum,
    kvartily.
    
    \item \textbf{Grafická kontrola:} histogramy, boxploty, scatterploty, korelační grafy.
    
    \item \textbf{Chybějící hodnoty:} mnoho redukčních metod neumí přímo pracovat s missing values.
    Je nutné je odstranit nebo imputovat.
    
    \item \textbf{Odlehlé hodnoty:} outliery mohou výrazně ovlivnit kovarianční/korelační matici,
    a tím i PCA, faktorovou i diskriminační analýzu.
    
    \item \textbf{Normalizace/standardizace:} proměnné v různých měřítkách mohou vést k tomu, že
    proměnná s velkou variabilitou dominuje výsledku.
\end{itemize}

Standardizace proměnné:
\[
z_{ij}
=
\frac{x_{ij}-\bar{x}_j}{s_j}.
\]

Používá se hlavně tehdy, když jsou proměnné v různých jednotkách, například příjem v Kč,
věk v letech a počet nákupů v kusech.

\subsection*{Odlehlé hodnoty v EDA}

Odlehlé pozorování je hodnota, která se výrazně liší od většiny dat. Může jít o chybu měření,
chybu zápisu, ale také o skutečně extrémní a informativní pozorování.

\subsubsection*{Z-score metoda}

\[
z_i = \frac{x_i-\bar{x}}{s}.
\]

Za potenciální outlier se často považuje pozorování:
\[
|z_i|>3.
\]

Tato metoda je jednoduchá, ale předpokládá přibližně normální rozdělení a není robustní vůči extrémům.

\subsubsection*{IQR metoda a Tukeyho hranice}

Interkvartilové rozpětí:
\[
IQR = Q_3-Q_1.
\]

Mírné outliery:
\[
x_i < Q_1 - 1.5IQR
\quad \text{nebo} \quad
x_i > Q_3 + 1.5IQR.
\]

Extrémní outliery:
\[
x_i < Q_1 - 3IQR
\quad \text{nebo} \quad
x_i > Q_3 + 3IQR.
\]

Výhodou je větší robustnost než u z-score metody.

\subsubsection*{Mahalanobisova vzdálenost}

Pro vícerozměrná data:
\[
D_i^2
=
(\mathbf{x}_i-\bar{\mathbf{x}})^\top
S^{-1}
(\mathbf{x}_i-\bar{\mathbf{x}}),
\]
kde $S$ je kovarianční matice.

Mahalanobisova vzdálenost bere v úvahu:
\begin{itemize}
    \item vzdálenost od vícerozměrného průměru,
    \item variabilitu proměnných,
    \item korelace mezi proměnnými.
\end{itemize}

Je proto vhodnější pro vícerozměrné outliery než posuzování každé proměnné zvlášť.

\subsection*{Redukce rozměru dat}

Redukční metody se používají, když je dataset příliš složitý:
\begin{itemize}
    \item má mnoho proměnných,
    \item proměnné spolu silně korelují,
    \item data se obtížně vizualizují,
    \item chceme odstranit redundanci,
    \item chceme získat menší počet nových, lépe interpretovatelných dimenzí.
\end{itemize}

Rozlišujeme:
\begin{itemize}
    \item \textbf{redukci počtu proměnných} -- PCA, faktorová analýza,
    \item \textbf{redukci počtu pozorování} -- shluková analýza,
    \item \textbf{redukci kategoriálních vztahů} -- korespondenční analýza,
    \item \textbf{redukci s ohledem na známé skupiny} -- diskriminační analýza.
\end{itemize}

\subsection*{Analýza hlavních komponent}

Analýza hlavních komponent, PCA, je metoda pro redukci počtu kvantitativních proměnných.
Původní proměnné nahradí menším počtem nových proměnných, které se nazývají hlavní komponenty.

\subsubsection*{Základní idea PCA}

PCA hledá lineární kombinace původních proměnných:
\[
z_1 = a_{11}x_1+a_{12}x_2+\cdots+a_{1p}x_p,
\]
\[
z_2 = a_{21}x_1+a_{22}x_2+\cdots+a_{2p}x_p,
\]
atd.

Tyto nové proměnné mají dvě důležité vlastnosti:
\begin{itemize}
    \item jsou navzájem nekorelované,
    \item první komponenta vysvětluje co největší část variability dat,
    druhá komponenta co největší část zbývající variability atd.
\end{itemize}

\subsubsection*{Maticový zápis PCA}

Mějme datovou matici:
\[
X \in \mathbb{R}^{n \times p},
\]
kde $n$ je počet pozorování a $p$ počet proměnných.

Předpokládejme, že proměnné jsou centrované, případně standardizované.

Kovarianční matice:
\[
S = \frac{1}{n-1}X^\top X.
\]

Pokud pracujeme se standardizovanými proměnnými, pak $S$ odpovídá korelační matici.

PCA je založená na vlastních číslech a vlastních vektorech:
\[
S\mathbf{v}_j = \lambda_j \mathbf{v}_j.
\]

Kde:
\begin{itemize}
    \item $\lambda_j$ je vlastní číslo,
    \item $\mathbf{v}_j$ je vlastní vektor,
    \item vlastní vektor určuje směr komponenty,
    \item vlastní číslo říká, kolik variability daná komponenta vysvětluje.
\end{itemize}

První hlavní komponenta odpovídá největšímu vlastnímu číslu:
\[
\lambda_1 \geq \lambda_2 \geq \cdots \geq \lambda_p.
\]

Skóre hlavních komponent:
\[
Z = XV,
\]
kde $V$ je matice vlastních vektorů.

\subsubsection*{Vysvětlená variabilita}

Celková variabilita:
\[
\sum_{j=1}^p \lambda_j.
\]

Podíl variability vysvětlený $j$-tou komponentou:
\[
\frac{\lambda_j}{\sum_{\ell=1}^p \lambda_\ell}.
\]

Kumulativní vysvětlená variabilita prvních $m$ komponent:
\[
\frac{\sum_{j=1}^m \lambda_j}
{\sum_{\ell=1}^p \lambda_\ell}.
\]

\subsubsection*{Volba počtu komponent}

Počet komponent lze volit podle:
\begin{itemize}
    \item \textbf{kumulativní vysvětlené variability} -- např. vybereme tolik komponent,
    aby vysvětlily 70--90~\% variability,
    
    \item \textbf{scree plotu} -- hledáme bod zlomu, po kterém další komponenty přidávají
    už jen málo informace,
    
    \item \textbf{Kaiserova kritéria} -- u PCA z korelační matice ponecháme komponenty s
    \[
    \lambda_j>1,
    \]
    
    \item \textbf{interpretovatelnosti} -- komponenty musí dávat věcný smysl.
\end{itemize}

\subsubsection*{Interpretace komponent}

Komponenty interpretujeme podle vah původních proměnných ve vlastních vektorech.
Pokud má první komponenta velké kladné váhy u proměnných příjem, útrata a hodnota majetku,
můžeme ji interpretovat například jako \uv{ekonomickou sílu}.

Často se používají také \textbf{loadings}, tedy vztahy mezi původními proměnnými a komponentami:
\[
\ell_{jk}
=
\operatorname{cor}(x_j,z_k).
\]

Velké absolutní hodnoty loadings říkají, že daná proměnná silně souvisí s danou komponentou.

\subsubsection*{Rekonstrukce dat}

Pokud ponecháme jen prvních $m<p$ komponent, dostaneme redukovanou aproximaci dat:
\[
\hat X_m = Z_m V_m^\top.
\]

Čím větší je $m$, tím přesnější je rekonstrukce, ale tím menší je redukce.

\subsubsection*{Výhody a nevýhody PCA}

Výhody:
\begin{itemize}
    \item snižuje počet proměnných,
    \item odstraňuje multikolinearitu,
    \item umožňuje vizualizovat vícerozměrná data ve 2D nebo 3D,
    \item zachovává co nejvíce variability.
\end{itemize}

Nevýhody:
\begin{itemize}
    \item komponenty nemusí být snadno interpretovatelné,
    \item PCA je citlivá na škálování proměnných,
    \item PCA je citlivá na outliery,
    \item PCA maximalizuje variabilitu, ne nutně vztah k cílové proměnné.
\end{itemize}

\subsection*{Regrese hlavních komponent}

Regrese hlavních komponent, PCR, kombinuje PCA a lineární regresi.

Nejprve se z původních vysvětlujících proměnných vytvoří hlavní komponenty:
\[
Z = XV.
\]

Potom se odhaduje regresní model:
\[
y = \alpha+\gamma_1 z_1+\cdots+\gamma_m z_m+u.
\]

PCR se používá hlavně při:
\begin{itemize}
    \item silné multikolinearitě vysvětlujících proměnných,
    \item velkém počtu proměnných,
    \item predikčních úlohách.
\end{itemize}

Nevýhoda je, že PCA vybírá komponenty podle variability v $X$, nikoli podle toho, jak dobře vysvětlují $y$.
Komponenta s malou variabilitou může být pro predikci $y$ důležitá, ale PCA ji může odstranit.

\subsection*{Faktorová analýza}

Faktorová analýza se podobá PCA, ale její cíl je jiný. PCA hledá nové proměnné, které vysvětlují
variabilitu dat. Faktorová analýza předpokládá, že za pozorovanými proměnnými stojí menší počet
nepozorovaných latentních faktorů.

\subsubsection*{Základní idea}

Příklad:
\begin{itemize}
    \item otázky v dotazníku o spokojenosti,
    \item ekonomické ukazatele regionů,
    \item psychologické škály,
    \item marketingové proměnné o zákaznících.
\end{itemize}

Mnoho pozorovaných proměnných může být projevem několika latentních faktorů, například:
\begin{itemize}
    \item socioekonomický status,
    \item loajalita zákazníka,
    \item kvalita života,
    \item finanční stabilita.
\end{itemize}

\subsubsection*{Faktorový model}

Standardní faktorový model:
\[
\mathbf{x}
=
\boldsymbol{\mu}
+
\Lambda \mathbf{f}
+
\boldsymbol{\varepsilon}.
\]

Kde:
\begin{itemize}
    \item $\mathbf{x}$ je vektor pozorovaných proměnných,
    \item $\boldsymbol{\mu}$ je vektor středních hodnot,
    \item $\mathbf{f}$ je vektor latentních faktorů,
    \item $\Lambda$ je matice faktorových zátěží,
    \item $\boldsymbol{\varepsilon}$ jsou specifické chyby.
\end{itemize}

Typicky předpokládáme:
\[
\mathrm{E}(\mathbf{f})=0,
\qquad
\mathrm{Var}(\mathbf{f})=I,
\qquad
\mathrm{E}(\boldsymbol{\varepsilon})=0,
\qquad
\mathrm{Cov}(\mathbf{f},\boldsymbol{\varepsilon})=0.
\]

Kovarianční matice pozorovaných proměnných:
\[
\Sigma
=
\Lambda\Lambda^\top+\Psi,
\]
kde $\Psi$ je diagonální matice unikátních rozptylů.

\subsubsection*{Faktorové zátěže}

Prvek $\lambda_{jk}$ v matici $\Lambda$ říká, jak silně je proměnná $x_j$ spojena s faktorem $f_k$.

Vysoká absolutní hodnota:
\[
|\lambda_{jk}|
\]
znamená, že proměnná silně souvisí s daným faktorem.

\subsubsection*{Komunalita a unikátnost}

Komunalita proměnné $x_j$:
\[
h_j^2
=
\sum_{k=1}^m \lambda_{jk}^2.
\]

Komunalita říká, jak velká část variability proměnné je vysvětlena společnými faktory.

Unikátnost:
\[
\psi_j = 1-h_j^2
\]
u standardizovaných proměnných.

Unikátnost vyjadřuje část variability, kterou společné faktory nevysvětlují.

\subsubsection*{Postup faktorové analýzy}

\begin{enumerate}
    \item \textbf{Příprava dat:}
    kontrola missing values, outlierů, škálování a korelací mezi proměnnými.
    
    \item \textbf{Posouzení vhodnosti dat:}
    proměnné by spolu měly dostatečně korelovat. Často se používá KMO kritérium a Bartlettův test
    sphericity.
    
    \item \textbf{Volba počtu faktorů:}
    scree plot, vlastní čísla, teorie, interpretovatelnost.
    
    \item \textbf{Extrakce faktorů:}
    například principal axis factoring nebo maximum likelihood.
    
    \item \textbf{Rotace faktorů:}
    rotace zjednodušuje interpretaci faktorové struktury.
    
    \item \textbf{Interpretace faktorů:}
    faktory pojmenujeme podle proměnných s nejvyššími zátěžemi.
\end{enumerate}

\subsubsection*{Rotace faktorů}

Rotace nemění základní kvalitu vysvětlení, ale snaží se dosáhnout jednodušší interpretace.

\paragraph{Ortogonální rotace.}
Faktory zůstávají nekorelované.
\begin{itemize}
    \item \textbf{Varimax} -- snaží se snížit křížové zátěže a zpřehlednit faktorový model.
    \item \textbf{Quartimax} -- snaží se snížit počet faktorů potřebných k vysvětlení proměnné.
    \item \textbf{Equamax} -- kompromis mezi varimax a quartimax.
\end{itemize}

\paragraph{Šikmá rotace.}
Faktory mohou být korelované.
\begin{itemize}
    \item \textbf{Promax} -- vhodná a efektivní i pro větší datasety.
    \item \textbf{Direct oblimin} -- může vést k jednodušší faktorové struktuře.
\end{itemize}

\subsubsection*{PCA vs. faktorová analýza}

\begin{itemize}
    \item \textbf{PCA} vytváří komponenty jako lineární kombinace pozorovaných proměnných.
    \item \textbf{Faktorová analýza} předpokládá latentní faktory, které způsobují korelace mezi
    pozorovanými proměnnými.
    \item PCA se zaměřuje na vysvětlení celkové variability.
    \item Faktorová analýza se zaměřuje na vysvětlení společné variability.
    \item PCA nemá explicitní chybový člen pro každou proměnnou.
    \item Faktorová analýza rozlišuje společnou a unikátní část variability.
\end{itemize}

\subsection*{Korespondenční analýza}

Korespondenční analýza je metoda pro exploraci vztahů v kontingenční tabulce. Používá se hlavně
pro kategoriální data.

\subsubsection*{Vstupní data}

Máme kontingenční tabulku:
\[
N = (n_{ij}),
\]
kde $n_{ij}$ je četnost v řádku $i$ a sloupci $j$.

Celkový počet pozorování:
\[
n = \sum_i \sum_j n_{ij}.
\]

Matice relativních četností:
\[
P = \frac{1}{n}N.
\]

Řádkové marginální četnosti:
\[
r_i = \sum_j p_{ij}.
\]

Sloupcové marginální četnosti:
\[
c_j = \sum_i p_{ij}.
\]

\subsubsection*{Princip}

Korespondenční analýza zkoumá, jak moc se skutečné četnosti liší od četností očekávaných při nezávislosti
řádkové a sloupcové proměnné.

Očekávaná relativní četnost při nezávislosti:
\[
p_{ij}^{(0)}=r_i c_j.
\]

Odchylka od nezávislosti:
\[
p_{ij}-r_i c_j.
\]

Testová statistika chí-kvadrát:
\[
\chi^2
=
\sum_i\sum_j
\frac{(n_{ij}-nr_ic_j)^2}{nr_ic_j}.
\]

Celková inertia:
\[
I
=
\frac{\chi^2}{n}.
\]

Inertia je analogie vysvětlené variability v PCA. Říká, jak silná je celková asociace mezi řádky a sloupci.

\subsubsection*{Standardizovaná reziduální matice}

Definujeme diagonální matice:
\[
D_r=\operatorname{diag}(r_1,\ldots,r_I),
\qquad
D_c=\operatorname{diag}(c_1,\ldots,c_J).
\]

Standardizovaná reziduální matice:
\[
S
=
D_r^{-1/2}
(P-\mathbf{r}\mathbf{c}^\top)
D_c^{-1/2}.
\]

Na tuto matici se aplikuje singulární rozklad:
\[
S
=
U\Delta V^\top.
\]

Diagonální prvky matice $\Delta$ jsou singulární hodnoty:
\[
\delta_1,\delta_2,\ldots.
\]

Inertia vysvětlená dimenzí $k$:
\[
\delta_k^2.
\]

Podíl vysvětlené inercie:
\[
\frac{\delta_k^2}{\sum_\ell \delta_\ell^2}.
\]

Maximální počet dimenzí:
\[
\min(I-1,J-1).
\]

\subsubsection*{Korespondenční mapa}

Řádkové souřadnice:
\[
F
=
D_r^{-1/2}U\Delta.
\]

Sloupcové souřadnice:
\[
G
=
D_c^{-1/2}V\Delta.
\]

V praxi se většinou zobrazují první dvě dimenze.

\subsubsection*{Interpretace}

\begin{itemize}
    \item Body řádků blízko sebe mají podobné profily přes sloupce.
    \item Body sloupců blízko sebe mají podobné profily přes řádky.
    \item Řádky a sloupce vzdálené od středu se více podílejí na závislosti.
    \item Body blízko středu odpovídají průměrnému profilu, tedy nejsou pro asociaci příliš specifické.
    \item Blízkost bodu řádku a bodu sloupce se interpretuje opatrně; nejčistší interpretace je
    uvnitř stejného typu bodů, tedy řádek--řádek a sloupec--sloupec.
\end{itemize}

\subsubsection*{Použití}

Korespondenční analýza se hodí například pro:
\begin{itemize}
    \item vztah značek a zákaznických segmentů,
    \item vztah politických stran a demografických skupin,
    \item vztah kategorií odpovědí v dotazníku,
    \item vizualizaci kontingenčních tabulek.
\end{itemize}

\subsection*{Diskriminační analýza}

Diskriminační analýza je metoda, která se používá, když známe skupiny pozorování a chceme najít kombinace
proměnných, které tyto skupiny co nejlépe oddělují.

Na rozdíl od PCA a faktorové analýzy je diskriminační analýza \textbf{supervised} metoda, protože používá
informaci o známé skupinové příslušnosti.

\subsubsection*{Základní situace}

Máme:
\begin{itemize}
    \item kvantitativní prediktory $x_1,\ldots,x_p$,
    \item kategoriální skupinovou proměnnou $g$,
    \item skupiny $1,\ldots,K$.
\end{itemize}

Cíl:
\begin{itemize}
    \item najít diskriminační funkce, které co nejlépe oddělují skupiny,
    \item klasifikovat nová pozorování do skupin,
    \item zjistit, které proměnné nejvíce přispívají k rozlišení skupin.
\end{itemize}

\subsubsection*{Fisherova diskriminační analýza}

Pro dvě skupiny hledáme lineární kombinaci:
\[
z_i = \mathbf{a}^\top \mathbf{x}_i,
\]
která maximalizuje rozdíl mezi skupinami vzhledem k variabilitě uvnitř skupin.

Fisherovo kritérium:
\[
J(\mathbf{a})
=
\frac{\mathbf{a}^\top S_B\mathbf{a}}
{\mathbf{a}^\top S_W\mathbf{a}},
\]
kde:
\begin{itemize}
    \item $S_B$ je meziskupinová variabilita,
    \item $S_W$ je vnitroskupinová variabilita.
\end{itemize}

Pro dvě skupiny:
\[
\mathbf{a}
\propto
S_W^{-1}(\bar{\mathbf{x}}_1-\bar{\mathbf{x}}_2).
\]

Tedy směr diskriminační funkce je dán rozdílem skupinových průměrů, upraveným o vnitroskupinovou
kovarianci.

\subsubsection*{Více skupin}

Pro více skupin definujeme:
\[
S_W
=
\sum_{k=1}^K
\sum_{i:g_i=k}
(\mathbf{x}_i-\bar{\mathbf{x}}_k)
(\mathbf{x}_i-\bar{\mathbf{x}}_k)^\top,
\]

\[
S_B
=
\sum_{k=1}^K
n_k
(\bar{\mathbf{x}}_k-\bar{\mathbf{x}})
(\bar{\mathbf{x}}_k-\bar{\mathbf{x}})^\top.
\]

Diskriminační směry řeší zobecněný vlastní problém:
\[
S_B\mathbf{a}
=
\lambda S_W\mathbf{a}.
\]

Maximální počet diskriminačních funkcí je:
\[
\min(p,K-1).
\]

\subsubsection*{Lineární diskriminační analýza}

Lineární diskriminační analýza, LDA, předpokládá, že skupiny mají stejnou kovarianční matici:
\[
\Sigma_1=\Sigma_2=\cdots=\Sigma_K=\Sigma.
\]

Klasifikační skóre pro skupinu $k$:
\[
\delta_k(\mathbf{x})
=
\mathbf{x}^\top \Sigma^{-1}\boldsymbol{\mu}_k
-
\frac{1}{2}
\boldsymbol{\mu}_k^\top
\Sigma^{-1}
\boldsymbol{\mu}_k
+
\log(\pi_k),
\]
kde:
\begin{itemize}
    \item $\boldsymbol{\mu}_k$ je vektor průměrů skupiny $k$,
    \item $\Sigma$ je společná kovarianční matice,
    \item $\pi_k$ je apriorní pravděpodobnost skupiny.
\end{itemize}

Pozorování přiřadíme do skupiny s největším skóre:
\[
\hat g(\mathbf{x})
=
\arg\max_k \delta_k(\mathbf{x}).
\]

\subsubsection*{Kvadratická diskriminační analýza}

Pokud skupiny nemají stejnou kovarianční matici, lze použít QDA:
\[
\Sigma_k \neq \Sigma_\ell.
\]

Klasifikační hranice pak nejsou lineární, ale kvadratické.

QDA je flexibilnější než LDA, ale potřebuje více dat, protože odhaduje samostatnou kovarianční matici pro
každou skupinu.

\subsubsection*{Předpoklady diskriminační analýzy}

\begin{itemize}
    \item kvantitativní vysvětlující proměnné,
    \item skupiny jsou předem známé,
    \item pozorování jsou nezávislá,
    \item přibližná vícerozměrná normalita v každé skupině,
    \item u LDA stejná kovarianční matice ve skupinách,
    \item bez silné multikolinearity mezi proměnnými.
\end{itemize}

\subsubsection*{Hodnocení klasifikace}

Výsledek diskriminační analýzy hodnotíme pomocí:
\begin{itemize}
    \item klasifikační matice,
    \item celkové přesnosti,
    \item chybovosti,
    \item sensitivity a specificity u dvou tříd,
    \item cross-validace.
\end{itemize}

Klasifikační matice porovnává skutečné a predikované skupiny:
\[
\begin{array}{c|cc}
 & \hat g=1 & \hat g=2 \\
\hline
g=1 & \text{správně} & \text{chyba} \\
g=2 & \text{chyba} & \text{správně}
\end{array}
\]

\subsection*{Multidimenzionální škálování}

Multidimenzionální škálování, MDS, převádí objekty do nízkorozměrného prostoru tak, aby byly co nejlépe
zachovány původní vzdálenosti mezi objekty.

MDS se hodí, když máme matici vzdáleností:
\[
D=(d_{ij}),
\]
ale ne nutně přímo původní proměnné.

Cíl:
\[
d_{ij}
\approx
\|\mathbf{z}_i-\mathbf{z}_j\|,
\]
kde $\mathbf{z}_i$ jsou nové souřadnice objektů v nízké dimenzi.

Často se minimalizuje stress:
\[
Stress
=
\sqrt{
\frac{
\sum_{i<j}(d_{ij}-\hat d_{ij})^2
}{
\sum_{i<j}d_{ij}^2
}
}.
\]

Nižší stress znamená lepší zachování původních vzdáleností.

\subsection*{Shluková analýza jako redukce pozorování}

Shluková analýza redukuje složitost dat tím, že seskupuje podobná pozorování do shluků.
Shluky by měly být:
\begin{itemize}
    \item vnitřně homogenní,
    \item navzájem odlišné,
    \item pokud možno dobře oddělené.
\end{itemize}

U hierarchického shlukování je důležité, jak měříme vzdálenost mezi shluky.

Single linkage:
\[
D(C_g,C_h)
=
\min_{\mathbf{x}_i\in C_g,\mathbf{x}_j\in C_h}
D(\mathbf{x}_i,\mathbf{x}_j).
\]

Complete linkage:
\[
D(C_g,C_h)
=
\max_{\mathbf{x}_i\in C_g,\mathbf{x}_j\in C_h}
D(\mathbf{x}_i,\mathbf{x}_j).
\]

Average linkage:
\[
D(C_g,C_h)
=
\frac{1}{n_g n_h}
\sum_{\mathbf{x}_i\in C_g}
\sum_{\mathbf{x}_j\in C_h}
D(\mathbf{x}_i,\mathbf{x}_j).
\]

Výsledkem hierarchického shlukování je dendrogram.

U nehierarchického shlukování, například $k$-means, se předem zvolí počet shluků $k$.
Algoritmus iterativně:
\begin{enumerate}
    \item zvolí počáteční středy shluků,
    \item přiřadí každé pozorování k nejbližšímu středu,
    \item přepočítá středy,
    \item opakuje postup do stabilizace.
\end{enumerate}

\subsection*{Srovnání hlavních metod}

\begin{itemize}
    \item \textbf{PCA:}
    používá kvantitativní proměnné, vytváří nekorelované komponenty, maximalizuje vysvětlenou
    variabilitu.
    
    \item \textbf{Faktorová analýza:}
    používá kvantitativní proměnné, hledá latentní faktory, vysvětluje korelační strukturu.
    
    \item \textbf{Korespondenční analýza:}
    používá kontingenční tabulku, analyzuje asociaci kategorií, je analogií PCA pro kategoriální data.
    
    \item \textbf{Diskriminační analýza:}
    používá známé skupiny, hledá funkce oddělující skupiny, slouží také ke klasifikaci.
    
    \item \textbf{MDS:}
    vychází z matice vzdáleností, snaží se zachovat vzdálenosti v nízkodimenzionálním zobrazení.
    
    \item \textbf{Shluková analýza:}
    redukuje počet objektů tím, že je seskupuje do podobných skupin.
\end{itemize}

\subsection*{Typická interpretace výstupů}

\begin{itemize}
    \item U PCA sledujeme vlastní čísla, vysvětlenou variabilitu, loadings a skóre komponent.
    \item U faktorové analýzy sledujeme faktorové zátěže, komunalitu, unikátnost a smysluplnost faktorů.
    \item U korespondenční analýzy sledujeme inercie a mapu kategorií.
    \item U diskriminační analýzy sledujeme diskriminační funkce, separaci skupin a klasifikační úspěšnost.
    \item U shlukování sledujeme vzdálenosti, dendrogram, počet shluků a interpretaci shluků.
\end{itemize}

\subsection*{Časté chyby v interpretaci}

\begin{itemize}
    \item Použití PCA bez standardizace, i když jsou proměnné v různých jednotkách.
    \item Interpretace komponent pouze podle názvu metody, nikoli podle loadings.
    \item Příliš mnoho nebo příliš málo ponechaných komponent/faktorů.
    \item Ignorování outlierů, které mohou výrazně změnit výsledek.
    \item Zaměňování PCA a faktorové analýzy.
    \item U korespondenční analýzy příliš silná interpretace vzdáleností mezi řádkovými a sloupcovými body.
    \item Použití diskriminační analýzy bez kontroly předpokladů a bez validace klasifikace.
\end{itemize}

\subsection*{Shrnutí otázky}

Exploratorní analýza slouží k pochopení dat, nalezení vzorů, odhalení problémů a přípravě dat pro další
modelování. Redukční metody pomáhají zjednodušit složitá data.

PCA nahrazuje původní kvantitativní proměnné menším počtem nekorelovaných hlavních komponent:
\[
S\mathbf{v}_j=\lambda_j\mathbf{v}_j.
\]
Komponenty jsou seřazeny podle vysvětlené variability.

Faktorová analýza předpokládá latentní faktory:
\[
\mathbf{x}
=
\boldsymbol{\mu}
+
\Lambda \mathbf{f}
+
\boldsymbol{\varepsilon},
\]
a snaží se vysvětlit korelační strukturu proměnných pomocí menšího počtu faktorů.

Korespondenční analýza analyzuje kontingenční tabulky. Vychází z odchylek skutečných četností od četností
očekávaných při nezávislosti a pracuje s inercií:
\[
I=\frac{\chi^2}{n}.
\]

Diskriminační analýza je supervised metoda. Hledá lineární kombinace proměnných, které co nejlépe oddělují
předem známé skupiny:
\[
J(\mathbf{a})
=
\frac{\mathbf{a}^\top S_B\mathbf{a}}
{\mathbf{a}^\top S_W\mathbf{a}}.
\]

Všechny tyto metody jsou citlivé na kvalitu dat, chybějící hodnoty, outliery a škálování proměnných.

\newpage
